% !TEX root =  ../Dissertation.tex

\chapter{Literature Review}
User retention in gaming is a new topic of conversation in scholarship,\cite{Tseng2015} as retention was previously less critical to the gaming space.\cite{Petrovskaya2022} Now that the Free-to-Play (F2P) model is more prevalent \cite{Koskenvoima2015}, particularly in mobile gaming,\cite{Skorin-Kapov2022} having a user stay on longer increases their chances of making in-app purchases.\cite{Jang2021} However, these same concepts can be used to retain users and monetise them differently.\cite{Jung2022} The daily game genre of web applications tends to use advertisements to monetise their websites.\cite{Framed} Another way is through memberships to the website that can include benefits such as exclusive content, the ability to export and import your stats, and the ability to sync your stats for all of their other sites.\cite{Ko-fi}

Even with these monetisation approaches, user retention can still benefit this genre of web application.\cite{Ross2018} It is shown that a wider breadth of engagement with any content increases the chance of a user becoming a member.\cite{Yang2022} While the quality and trust of a website are not indicative of how many advertisements someone will click, as it is more based on the quality of the advert, the chance of a user clicking on an advert will naturally be linked with how much time a user spends on a website.\cite{Werner2022}

While it has been established that user retention is still important for web applications without micro-transactions, it is also necessary to decide what sort of user retention techniques are the strongest without adding unnecessary monetisation and keeping users satisfied.\cite{Ross2018} To find these techniques, World of Warcraft is a fitting start as they have been using a Games as a Service model (GaaS) for many years before the popularisation of F2P and mobile gaming. Thomas Debeauvais et al. have researched the retention mechanisms in World of Warcraft, and their research indicated two main techniques that are driving factors that retain players: achievement and social play.\cite{Debeauvais2010} The ‘hardcore’ players in World of Warcraft have higher playtime and, in surveys, were more interested in their gaming experience, having a challenge and a reward for their achievements. Casual players are more interested in the game's social aspect, which involves players grouping with their friends to take on challenges together.\cite{Debeauvais2010}

When researchers explore mobile play, these two retention techniques of ‘social’ and ‘achievement’ are consistent. Research has been done into the MMO World of Tanks players, which are also available as an app on mobile phones. The game is centred around armoured warfare, where players control a variety of tanks and engage in team-based battles. The game emphasises strategy, teamwork, and tactical combat rather than fast-paced action. When Liyang Gu et al. researched the World of Tanks replay system,\cite{Gu2021}, they found that while only a small number of users used this system, they tended to be more ‘hardcore’ players. This is because the replay communities surrounding this system mean that users felt a social aspect of the game and were more dedicated to playing. Immersion is something else in World of Tanks that affects user retention. A different paper by Liyang Gu et al. shows that users with a stronger sense of immersion in World of Tanks are also more likely to be retained.\cite{Gu2018} 

Paul Bertens et al. have made a scalable multidimensional churn prediction model that has been very accurate at predicting player churn.\cite{Bertens2017} One of the main input predictors for player churn is that the higher the level of the player, the less likely they are to churn. This indicates that the more a user feels they have achieved in a game, the more likely they will be retained for extended periods.\cite{Bertens2017}

Another retention technique is adding features, which can improve the number of users. However, it is essential to note that adding too many features can affect user acquisition, as a user may try to access your product and find it ‘overwhelming’ or ‘confusing,’ which may lead them to use another product.\cite{Hamilton2017}


\section{Ethics of Retention Techniques}
Various retention techniques have been studied and applied to online games. Some games have used these techniques, especially the more ethically dubious ones, to push their game to levels of addiction where players will find themselves spending an excessive amount of time on a game and often spending an excessive amount of money. A leading force behind this addiction that players can find themselves in is the retention techniques that have been applied to these games.\cite{Bernik2023}

A genre of game that has mastered retention techniques is ‘gacha games’. Gacha games are a genre primarily founded in Japan in the early 2010s. The etymology of the word ‘Gacha’ comes from the word ‘Gachapon’; these vending machines originated in Japan in the 1960s and dispensed different toys each time you paid.\cite{Bernik2023} Gacha games have continued that legacy by having in-game rewards that depend on random chance, which can be increased with in-app purchases of ‘tickets’ that will let you roll multiple times for a reward.\cite{Bernik2023}

Many have criticised Gacha games for their use of retention techniques as they bring gambling into the hands of younger players. While the ethics of this gameplay loop are questionable, its efficacy cannot be denied. miHoyo, one of the largest gacha game studios, has gone from a company with 15 million dollars of revenue in 2015 to 1.7 billion dollars in 2023.\cite{Pingwest2023} A large part of this is their use of the ‘free to play (F2P)’ and ‘games as a service (GaaS)’ models that heavily rely on retention techniques to keep their player base consistently returning to their game.\cite{Bernik2023}

Genshin Impact was one of MiHoyo's biggest games, recently overtaken by more recent titles but still making 60,500,000 USD in revenue in June 2024 and the game that introduced them to Western audiences.\cite{Gacha2024} It applies all the previously mentioned techniques: regular content updates, a gacha system, daily and weekly activities, wait-to-play, social features, progression, and reward systems. This led to massive success with both Asian and Western audiences.\cite{Greting2022}

Recent scholarship has pushed back against the use of these techniques. One of whom is Harry Brignull, a doctor in Cognitive Science who coined the phrase ‘dark patterns’ in his book "Deceptive Patterns: Exposing the Tricks Tech Companies Use to Control You" \cite{Brignull2023}. This phrase is often used in online discourse in video games like Genshin Impact. Harry Brignull argues that retention techniques like waiting to play, daily activities, and invested value intentionally manipulate cognitive biases to create an addiction and exploit the user for monetary gain.\cite{Brignull2023}

While this paper intends to advocate for the use of some retention techniques, it is essential to point out that the intention is not to create a version of Wordle that uses these ‘dark pattern’ retention techniques to the end of exploiting the user by intentionally making it too addictive nor should this be construed as an advocacy of the inclusion of microtransactions in the daily game genre of web applications. This paper advocates for a middle ground of adding proven retention techniques that increase user satisfaction and retention without adding these dark pattern techniques used in the gacha game genre to exploit the user for monetary gain. 
